\documentclass[11pt,letterpaper]{article}

% No additional packages are required.
% Compile with: pdflatex assignment1.tex (run twice for the references).
\pdfpagewidth=\paperwidth
\pdfpageheight=\paperheight
\setlength{\oddsidemargin}{0pt}
\setlength{\evensidemargin}{0pt}
\setlength{\textwidth}{6.5in}
\setlength{\topmargin}{0pt}
\setlength{\headheight}{0pt}
\setlength{\headsep}{0pt}
\setlength{\textheight}{9in}
\setlength{\footskip}{0.4in}
\setlength{\parindent}{0pt}
\setlength{\parskip}{3pt}
\renewcommand{\labelenumi}{(\alph{enumi})}
\newcommand{\question}[1]{%
  \par\medskip\noindent{\large\bfseries Question #1}\par\nobreak\smallskip}
\frenchspacing

\begin{document}

\begin{center}
  {\Large\bfseries ECE 759: Assignment 1}\\[4pt]
  Xusheng Zhi (\texttt{xzhi8})
\end{center}

\textbf{GitHub (HW01):}
{\small\texttt{https://github.com/Hihi142/repo759/tree/main/HW01}}\\
\textbf{SSH:} \texttt{git@github.com:Hihi142/repo759.git}

\question{1}
I went through a) through c) and understand how to time code, how to submit
my assignments with git, and what the recommended workflow is when it comes
to working on my assignment.

\question{2}
\begin{tabular}{@{}ll@{}}
(a) & \verb|cd somedir|\\[3pt]
(b) & \verb|cat sometext.txt|\\[3pt]
(c) & \verb|tail -n 5 sometext.txt|\\[3pt]
(d) & \verb|tail -n 5 ./*.txt|\\[3pt]
(e) & \verb|for i in {0..6}; do echo "$i"; done|
\end{tabular}

\question{3}
\begin{enumerate}
\setlength{\itemsep}{3pt}
\setlength{\parsep}{0pt}
\item No modules are loaded at login; \texttt{module list} reports
\texttt{No modules loaded}.

\item The default GCC version is \texttt{14.3.1}
(\texttt{Red Hat 14.3.1-4}).

\item The available CUDA modules are:
\begin{center}
\small
\begin{tabular}{@{}ll@{}}
\texttt{nvidia/cuda/10.2.2} & \texttt{nvidia/cuda/11.0.3}\\
\texttt{nvidia/cuda/11.3.1} & \texttt{nvidia/cuda/11.6.0}\\
\texttt{nvidia/cuda/11.8.0} & \texttt{nvidia/cuda/12.0.0}\\
\texttt{nvidia/cuda/12.1.0} & \texttt{nvidia/cuda/12.2.0}\\
\texttt{nvidia/cuda/12.5.0} & \texttt{nvidia/cuda/12.9.1}\\
\texttt{nvidia/cuda/13.0.0} (default) &
\end{tabular}
\end{center}
The CUDA search also returned these related modules:
\begin{center}
\small
\begin{tabular}{@{}l@{}}
\texttt{nvidia/nvhpc-hpcx-cuda11/24.5}\\
\texttt{nvidia/nvhpc-hpcx-cuda12/23.11}\\
\texttt{nvidia/nvhpc-hpcx-cuda12/24.5} (default)
\end{tabular}
\end{center}

\item \texttt{valgrind/3.26.0} is a dynamic analysis tool used for
memory debugging.
\end{enumerate}

\question{4}
Source: \texttt{HW01/task4.sh}; submitted on Euler with
\texttt{sbatch task4.sh}.\\
Job name: \texttt{FirstSlurm}; requested resources: one task, two CPU cores.

\begin{tabular}{@{}ll@{}}
\texttt{FirstSlurm.out}: & \texttt{euler01.engr.wisc.edu}\\
\texttt{FirstSlurm.err}: & empty (no error output).
\end{tabular}

\newpage

\question{5}
\begin{enumerate}
\setlength{\itemsep}{5pt}
\setlength{\parsep}{0pt}
\item By default, a Slurm job starts in the directory where
\texttt{sbatch} was invoked, unless \texttt{--chdir} specifies
another directory.~\cite{sbatch}

\item \verb|SLURM_JOB_ID| is the environment variable containing
the current job ID.~\cite{sbatch}

\item I would use \texttt{squeue -u xzhi8} to check the status of my
pending, running, and completing jobs.~\cite{squeue}

\item I would use \verb|scancel JOB_ID|, replacing \verb|JOB_ID|
with the ID of the queued job to cancel.~\cite{scancel}

\item \texttt{\#SBATCH --gres=gpu:1} requests one GPU on each node
allocated to the job.~\cite{sbatch}

\item \textbf{(Optional)}
\texttt{\#SBATCH --array=0-9} creates a job array of ten tasks indexed
0 through 9, with each task's index available in
\verb|SLURM_ARRAY_TASK_ID|.~\cite{arrays}
\end{enumerate}

\question{6}
Source: \texttt{HW01/task6.cpp}.
The program reads a nonnegative integer \(N\), prints 0 through \(N\)
using \texttt{printf}, and prints \(N\) through 0 using \texttt{std::cout},
with space-separated integers and a newline after each sequence.

Compilation command:
\begin{verbatim}
g++ task6.cpp -Wall -O3 -std=c++17 -o task6
\end{verbatim}
Local test output for \texttt{./task6 6}:
\begin{verbatim}
0 1 2 3 4 5 6
6 5 4 3 2 1 0
\end{verbatim}

\begin{small}
\begin{thebibliography}{4}
\setlength{\itemsep}{2pt}
\setlength{\parsep}{0pt}
\bibitem{sbatch} SchedMD, \emph{sbatch manual}.\\
\texttt{https://slurm.schedmd.com/sbatch.html}
\bibitem{squeue} SchedMD, \emph{squeue manual}.\\
\texttt{https://slurm.schedmd.com/squeue.html}
\bibitem{scancel} SchedMD, \emph{scancel manual}.\\
\texttt{https://slurm.schedmd.com/scancel.html}
\bibitem{arrays} SchedMD, \emph{Job Array Support}.\\
\verb|https://slurm.schedmd.com/job_array.html|
\end{thebibliography}
\end{small}

\end{document}
